\expandafter\ifx\csname CAPETINCLUDEMODE\endcsname\relax
\documentclass[11pt,a4paper]{article}
% preamble.tex — packages et commandes communes a toutes les fiches CAPET SII IE
\input{glyphtounicode}
\pdfgentounicode=1
\usepackage{cmap}
\usepackage[T1]{fontenc}
\usepackage[utf8]{inputenc}
\usepackage[french]{babel}
\usepackage[a4paper, top=1.6cm, bottom=1.6cm, left=1.4cm, right=1.4cm]{geometry}
\setlength{\headheight}{14.5pt}
\usepackage{amsmath, amssymb, amsfonts}
\usepackage{siunitx}
\sisetup{locale=FR, output-decimal-marker={,}, inter-unit-product=\ensuremath{{}\cdot{}}}
\usepackage{array, booktabs, multirow, tabularx}
\usepackage{xcolor}
\usepackage[most, breakable]{tcolorbox}
\usepackage{multicol}
\usepackage{enumitem}
\usepackage{mathtools}
\usepackage{lmodern}
\usepackage{microtype}
\usepackage{fancyhdr}
\usepackage{tikz}
\usetikzlibrary{positioning,arrows.meta}
\usepackage{lastpage}

% ── Couleurs par domaine ──────────────────────────────────────────────
\definecolor{colElec}{RGB}{26, 82, 148}       % Électrotechnique — bleu foncé
\definecolor{colElecLight}{RGB}{210, 225, 245} % fond bleu clair
\definecolor{colAuto}{RGB}{34, 120, 66}        % Automatique — vert
\definecolor{colMeca}{RGB}{185, 90, 20}        % Mécanique — orange
\definecolor{colTherm}{RGB}{160, 40, 40}       % Thermique — rouge
\definecolor{colQT}{RGB}{232, 244, 233}        % fond question type (vert très léger)
\definecolor{colJury}{RGB}{255, 243, 205}      % fond jury (jaune ambre)
\definecolor{colJuryFrame}{RGB}{200, 130, 0}   % bordure jury
\definecolor{colGrayLight}{RGB}{245, 245, 247} % fond gris neutre
\definecolor{colSeparator}{RGB}{180, 180, 180} % gris séparateur

% ── Style tcolorbox de section principale ────────────────────────────
\tcbset{
  ficheSection/.style={
    breakable,
    enhanced,
    colback=#1!7!white,
    colframe=#1,
    colbacktitle=#1,
    coltitle=white,
    fonttitle=\bfseries\normalsize,
    left=4pt, right=4pt, top=3pt, bottom=4pt,
    before skip=6pt, after skip=6pt,
  }
}

% ── Style encadré question type ──────────────────────────────────────
\tcbset{
  questiontype/.style={
    breakable,
    enhanced,
    colback=colQT,
    colframe=colElec!55!black,
    colbacktitle=colElec!55!black,
    coltitle=white,
    fonttitle=\bfseries\small,
    left=4pt, right=4pt, top=3pt, bottom=3pt,
    before skip=5pt, after skip=5pt,
  }
}

% ── Style encadré jury ───────────────────────────────────────────────
\tcbset{
  jury/.style={
    breakable,
    enhanced,
    colback=colJury,
    colframe=colJuryFrame,
    fonttitle=\bfseries\normalsize,
    left=4pt, right=4pt, top=3pt, bottom=3pt,
  }
}

% ── Mise en page ─────────────────────────────────────────────────────
\pagestyle{fancy}
\fancyhf{}
\renewcommand{\headrulewidth}{0.4pt}
\renewcommand{\footrulewidth}{0.3pt}

\setlist[itemize]{leftmargin=1.2em, itemsep=1pt, topsep=2pt, parsep=0pt}
\setlist[enumerate]{leftmargin=1.2em, itemsep=3pt, topsep=2pt, parsep=0pt}

\setlength{\parskip}{2pt}
\setlength{\columnsep}{10pt}

\def\CAPETBEGINFILE{\begin{document}}
\def\CAPETENDFILE{\end{document}}
\else
\def\CAPETBEGINFILE{}
\def\CAPETENDFILE{}
\fi

\fancyhead[L]{\textbf{\textcolor{colElec}{T13}} - Machine asynchrone (MAS)}
\fancyhead[R]{\small\textit{CAPET SII IE - Session 2026}}
\fancyfoot[C]{\small\thepage\,/\,\pageref{LastPage}}

\CAPETBEGINFILE

% -----------------------------------------------------------------------
%  BLOC EN-TeTE
% -----------------------------------------------------------------------

\begin{tcolorbox}[
  enhanced, breakable,
  colback=colElec,
  colframe=colElec,
  left=6pt, right=6pt, top=5pt, bottom=5pt,
  before skip=0pt, after skip=4pt
]
  {\color{white}\bfseries\large T13 \quad Machine asynchrone (MAS)}\\[2pt]
  {\color{colElec!30!white}\small
    \textbf{Domaine :} electrotechnique - Machines electriques \quad
    \textbf{Presence :} \textbf{9/13} sujets (2013-2026) - theme \emph{tres frequent}
  }\\[1pt]
  {\color{colElec!20!white}\footnotesize
    Sous-themes : vitesse synchrone et glissement - bilan de puissance (Pa - Pt - Pmeca - Pu) - couple electromagnetique - schema equivalent monophase - demarrage etoile-triangle - commande V/f (voir T14)
  }
\end{tcolorbox}

% -----------------------------------------------------------------------
%  1. RAPPELS DE COURS
% -----------------------------------------------------------------------

\begin{tcolorbox}[
  ficheSection=colElec,
  title={1. Rappels de cours},
]

% -- 1.1 Vitesse synchrone et glissement ------------------------------
\textbf{1.1 Vitesse de synchronisme et glissement}

\smallskip
\begin{center}
\renewcommand{\arraystretch}{1.5}
\begin{tabular}{lll}
\toprule
\textbf{Grandeur} & \textbf{Expression} & \textbf{Unite} \\
\midrule
Vitesse synchrone & $N_s = \dfrac{60\,f}{p}$ & tremin$^{-1}$ \\[6pt]
Glissement & $g = \dfrac{N_s - N}{N_s} = \dfrac{\omega_s - \omega}{\omega_s}$ & - \\[6pt]
Vitesse pulsation rotorique & $\omega_r = g\,\omega_s$ & rades$^{-1}$ \\
\bottomrule
\end{tabular}
\end{center}
\textit{Methode : calculer $p = 60f/N_{nom}$, arrondir - l'entier le plus proche, puis recalculer $N_s$.}

\medskip\hrule\vspace{4pt}

% -- 1.2 Bilan de puissance --------------------------------------------
\textbf{1.2 Bilan de puissance - flux de droite - gauche}

\smallskip
\[
  P_a \xrightarrow{-P_{Js}-P_{Fs}} P_t \xrightarrow{-P_{Jr} = -g P_t} P_{meca} = (1-g)P_t \xrightarrow{-P_{pm}} P_u = C_u\Omega
\]

\begin{center}
\renewcommand{\arraystretch}{1.4}
\begin{tabular}{lll}
\toprule
\textbf{Puissance} & \textbf{Expression} & \textbf{Remarque} \\
\midrule
Absorbee au stator & $P_a = 3\,V_s I_s \cos\varphi$ & $V_s$ = tension \emph{simple} \\
Pertes Joule stator & $P_{Js} = 3\,R_s I_s^2$ & souvent negligees \\
Transmise au rotor & $P_t = P_a - P_{Js} - P_{Fs}$ & puissance de l'entrefer \\
Pertes Joule rotor & $P_{Jr} = g\,P_t = 3 R'_R I'^2_r$ & \\
Puissance mecanique & $P_{meca} = (1-g)\,P_t$ & developpee au rotor \\
Puissance utile & $P_u = P_{meca} - P_{pm} = C_u\,\Omega$ & sortie d'arbre \\
\bottomrule
\end{tabular}
\end{center}

\medskip\hrule\vspace{4pt}

% -- 1.3 Couple electromagnetique -------------------------------------
\textbf{1.3 Couple electromagnetique}

\smallskip
\[
  C_{em} = \frac{P_t}{\Omega_s}
  \qquad;\qquad
  C_u = C_{em} - C_p
  \qquad;\qquad
  C_{em} \approx \frac{3p\,V_s^2}{R'_R\,\omega_s^2} \cdot g \quad \text{(glissement faible)}
\]
Forme linearisee : $C_{em} = A \cdot \omega_r$ avec $A = \dfrac{3p V_s^2}{R'_R \omega_s^2}$ - valable pour $g \ll 1$.

\medskip\hrule\vspace{4pt}

% -- 1.4 Modele de Laplace et demarrage -------------------------------
\textbf{1.4 equation mecanique et demarrage etoile-triangle}

\smallskip
equation mecanique dans le domaine de Laplace :
\[
  \Omega(p) = \frac{K_{meca}}{1 + \tau_{meca}\,p}\,\bigl[C_{em}(p) - C_r(p)\bigr]
  \qquad
  K_{meca} = \frac{1}{f}\,,\quad \tau_{meca} = \frac{J}{f}
\]

\textbf{Demarrage etoile-triangle :} en etoile, chaque enroulement reeoit $V_{reseau}/\sqrt{3}$ - courant de demarrage divise par 3 - couple de demarrage divise par 3. Passage en triangle - quasi-vitesse nominale.

\end{tcolorbox}

% -----------------------------------------------------------------------
%  2. FORMULAIRE SYNTHeTIQUE
% -----------------------------------------------------------------------

\begin{tcolorbox}[
  ficheSection=colElec,
  title={2. Formulaire synthetique},
]
\small
\begin{multicols}{2}
\renewcommand{\arraystretch}{1.45}
\begin{tabular}{@{}p{4.5cm}ll@{}}
\toprule
\textbf{Formule} & \textbf{Usage} & \textbf{U.} \\
\midrule
\multicolumn{3}{@{}l}{\footnotesize\textit{Vitesse et glissement}} \\
$N_s = 60f/p$ & Synchro & tr/min \\
$g = (N_s-N)/N_s$ & Glissement & - \\
$\omega_r = g\,\omega_s$ & Puls. rotorique & rad/s \\
\midrule
\multicolumn{3}{@{}l}{\footnotesize\textit{Puissances}} \\
$P_a = 3\,V_s I_s\cos\varphi$ & Absorbee & W \\
$P_{Jr} = g\,P_t$ & Joule rotor & W \\
$P_{meca} = (1-g)\,P_t$ & Mecanique & W \\
$P_u = P_{meca} - P_{pm}$ & Utile & W \\
\bottomrule
\end{tabular}

\columnbreak

\renewcommand{\arraystretch}{1.45}
\begin{tabular}{@{}p{4.5cm}ll@{}}
\toprule
\textbf{Formule} & \textbf{Usage} & \textbf{U.} \\
\midrule
\multicolumn{3}{@{}l}{\footnotesize\textit{Couple}} \\
$C_{em} = P_t/\Omega_s$ & Couple em & Nem \\
$C_u = C_{em} - C_p$ & Couple utile & Nem \\
$C_{em} \approx A\cdot\omega_r$ & Linearise & Nem \\
\midrule
\multicolumn{3}{@{}l}{\footnotesize\textit{Modele Laplace}} \\
$K_{meca} = 1/f$ & Gain mec. & - \\
$\tau_{meca} = J/f$ & Constante & s \\
\midrule
\multicolumn{3}{@{}l}{\footnotesize\textit{Donnees 2024}} \\
$p = 2$, $N_s = 1500$ tr/min & Extracteur & - \\
$g = 6{,}67\%$, $P_u = 100$ kW & Extracteur & - \\
\bottomrule
\end{tabular}
\end{multicols}
\normalsize
\end{tcolorbox}

% -----------------------------------------------------------------------
%  3. QUESTIONS TYPES
% -----------------------------------------------------------------------

\begin{tcolorbox}[
  ficheSection=colElec,
  title={3. Questions types - extraites des sujets reels},
]

% -- QT1 -------------------------------------------------------------

% Mini-procedure de resolution type copie concours
\begin{tcolorbox}[questiontype,title={Procedure type copie concours (5 etapes)}]
\begin{enumerate}[leftmargin=1.5em]
  \item \textbf{Qualifier le modele} : hypotheses, domaine de validite, grandeurs d'entree/sortie.
  \item \textbf{Ecrire les lois} : equations de base avec symboles et unites SI.
  \item \textbf{Isoler l'inconnue} : expression litterale, puis application numerique tracee.
  \item \textbf{Controler l'ordre de grandeur} : coherence physique et unite finale.
  \item \textbf{Conclure concours} : phrase courte orientee decision technique.
\end{enumerate}
\end{tcolorbox}

% Visuel de methode : chemin de resolution CAPET
\begin{center}
\begin{tikzpicture}[font=\small, node distance=7mm and 6mm, >=Stealth]
  \node[draw, rounded corners, fill=colElec!8, inner sep=4pt] (n1) {Notions};
  \node[draw, rounded corners, fill=colElec!8, right=of n1, inner sep=4pt] (n2) {Methode};
  \node[draw, rounded corners, fill=colElec!8, right=of n2, inner sep=4pt] (n3) {Application annale};
  \node[draw, rounded corners, fill=colElec!8, right=of n3, inner sep=4pt] (n4) {Verification jury};
  \draw[->, thick, colElec] (n1) -- (n2);
  \draw[->, thick, colElec] (n2) -- (n3);
  \draw[->, thick, colElec] (n3) -- (n4);
\end{tikzpicture}
\end{center}
\begin{tcolorbox}[questiontype,
  title={QT1 - Vitesse synchrone, glissement, puissance absorbee \hfill\normalfont\footnotesize E2 disc. 2024 (A1-A7, extracteur d'air)}
]
\textbf{enonce.}
Plaque signaletique du moto-extracteur : $V_{composee} = 400\,\mathrm{V}$, $V_s = 230\,\mathrm{V}$, $I_s = 229\,\mathrm{A}$, $\cos\varphi_{mot} = 0{,}85$, $f = 50\,\mathrm{Hz}$, $N = 1400\,\mathrm{tr/min}$, $p = 2$.
Calculer $N_s$, $g$ et $P_a$.

\tcbline

\textbf{Corrige.}
\[
  N_s = \frac{60 \times 50}{2} = \mathbf{1500\,\mathrm{tr/min}}
\]
\[
  g = \frac{N_s - N}{N_s} = \frac{1500 - 1400}{1500} = 0{,}0\overline{6} \approx \mathbf{6{,}67\,\%}
\]
\[
  P_a = 3\,V_s I_s \cos\varphi = 3 \times 230 \times 229 \times 0{,}85 \approx \mathbf{134{,}6\,\mathrm{kW}}
\]
\textit{Rappel : utiliser la tension \emph{simple} $V_s = 230\,\mathrm{V}$ (et non 400\,V) avec le facteur 3.}
\end{tcolorbox}

\vspace{4pt}

% -- QT2 -------------------------------------------------------------

\begin{tcolorbox}[questiontype,
  title={QT2 - Bilan de puissance complet \hfill\normalfont\footnotesize E2 disc. 2024 (A5-A7, extracteur)}
]
\textbf{enonce.}
Donnees DT1 : $R'_R = 0{,}1\,\Omega$, $I'_r = 122\,\mathrm{A}$, $g = 0{,}067$, $C_p = 200\,\mathrm{N\!\cdot\!m}$, $N = 1400\,\mathrm{tr/min}$.
Calculer $P_t$, $P_{Jr}$, $P_{meca}$ et $P_u$.

\tcbline

\textbf{Corrige.}
\[
  P_t = \frac{3 R'_R}{g} I'^2_r = \frac{3 \times 0{,}1}{0{,}067} \times 122^2 \approx 66{,}8\,\mathrm{kW}
\]
\[
  P_{Jr} = g\,P_t = 0{,}067 \times 66{,}8 \approx 4{,}5\,\mathrm{kW}
  \qquad;\qquad
  P_{meca} = (1-g)\,P_t = 0{,}933 \times 66{,}8 \approx 62{,}3\,\mathrm{kW}
\]
\[
  \Omega = \frac{1400 \times 2\pi}{60} = 146{,}6\,\mathrm{rad\!\cdot\!s^{-1}}
  \quad;\quad
  P_{pm} = C_p\,\Omega = 200 \times 146{,}6 = 29{,}3\,\mathrm{kW}
  \quad;\quad
  P_u = 62{,}3 - 29{,}3 = \mathbf{33\,\mathrm{kW}}
\]
\textit{Note : le sujet donne $P_u \approx 100\,\mathrm{kW}$ avec les vraies donnees du DT1. La methode est identique.}
\end{tcolorbox}

\vspace{4pt}

% -- QT3 -------------------------------------------------------------

\begin{tcolorbox}[questiontype,
  title={QT3 - Couple EM linearise et modele Laplace \hfill\normalfont\footnotesize E2 disc. 2024 (C1-C3)}
]
\textbf{enonce.}
Montrer que pour $g \approx 0$ : $C_{em} = \dfrac{3p\,V_s^2}{R'_R\,\omega_s^2}\cdot g$.
En deduire la fonction de transfert mecanique $\Omega(p)$ avec $J = 0{,}035\,\mathrm{kg\!\cdot\!m^2}$ et $f = 4 \times 10^{-3}\,\mathrm{N\!\cdot\!m\!\cdot\!s\!\cdot\!rad^{-1}}$.

\tcbline

\textbf{Corrige.}

Pour $g \ll 1$ : $I'_r \approx V_s g / R'_R$ et $P_t = 3(R'_R/g)I'^2_r = 3V_s^2 g/R'_R$
\[
  C_{em} = \frac{P_t}{\Omega_s} = \frac{3V_s^2 g}{R'_R \omega_s} = \underbrace{\frac{3p V_s^2}{R'_R \omega_s^2}}_{A} \cdot \underbrace{g\omega_s}_{\omega_r} = A\cdot\omega_r
\]

\textbf{Modele Laplace :} $J\,d\Omega/dt + f\,\Omega = C_{em} - C_r \;\Rightarrow\; \Omega(p) = \dfrac{K_{meca}}{1+\tau_{meca}\,p}\,[C_{em}(p)-C_r(p)]$
\[
  K_{meca} = \frac{1}{f} = \frac{1}{4\times10^{-3}} = 250\,\mathrm{rad\!\cdot\!s^{-1}\!\cdot\!N^{-1}\!\cdot\!m^{-1}}
  \qquad;\qquad
  \tau_{meca} = \frac{J}{f} = \frac{0{,}035}{0{,}004} = 8{,}75\,\mathrm{s}
\]
\end{tcolorbox}

\end{tcolorbox}

% -----------------------------------------------------------------------
%  4. ATTENTES DU JURY
% -----------------------------------------------------------------------

\begin{tcolorbox}[
  jury,
  title={4. Attentes du jury - Erreurs recurrentes et points de vigilance},
]

\begin{enumerate}[label=\textbf{\arabic*.}, leftmargin=*, itemsep=5pt]

  \item \textbf{Facteur 3 en triphase - erreur tres frequente.}
  $P_a = \mathbf{3}\,V_s I_s\cos\varphi$ et non $V_s I_s\cos\varphi$. Utiliser la tension \emph{simple} $V_s = 230\,\mathrm{V}$ (et non la composee 400\,V) si le facteur 3 est present. Rapport jury 2024 : lacune signalee.

  \item \textbf{$p$ entier obligatoire.}
  Calculer $p = 60f/N_{nom}$, puis arrondir - l'entier le plus proche, puis recalculer $N_s$. Le jury 2023 insiste : "verifier que le resultat est un entier".

  \item \textbf{Glissement en decimal dans les formules.}
  $g = 0{,}067$ (decimal) dans $P_{Jr} = g\,P_t$, pas $g = 6{,}67$.

  \item \textbf{$P_t \neq P_{meca}$.}
  $P_t$ est la puissance transmise - l'entrefer (puissance de l'entrefer). $P_{meca} = (1-g)P_t$ est la puissance mecanique developpee. $P_{Jr} = g\,P_t$ est perdue par effet Joule au rotor.

  \item \textbf{Demarrage etoile-triangle.}
  Le couple de demarrage est divise par 3 en etoile - demarrage - charge reduite uniquement. Ne pas conclure que c'est adapte - toutes les charges.

  \item \textbf{ecrire l'expression litterale puis l'application numerique.}
  Rapport jury 2024 : "les resultats donnes directement, sans calcul et sans justification ne repondent pas aux attendus". Trois lignes par calcul minimum.

\end{enumerate}

\medskip
\begin{tcolorbox}[
  colback=colJuryFrame!12!colJury,
  colframe=colJuryFrame,
  left=4pt, right=4pt, top=2pt, bottom=2pt,
  fontupper=\small\itshape
]
Rapport jury 2024 : "L'objectif de cette partie est de dimensionner le variateur et du moto-extracteur electrique pour optimiser la consommation energetique de l'extracteur d'air. Seulement 55 \% des candidats l'ont traitee correctement."\\
Rapport jury 2024 : "Le jury regrette le manque de maetrise des formules de base sur les reseaux triphases, le regime sinusoedal et les relations d'electrocinetique."\\
Rapport jury 2023 : "La majorite des candidats maetrise la formule de la vitesse de synchronisme mais beaucoup oublient de verifier que le resultat est un entier pour le nombre de paires de peles."
\end{tcolorbox}

\end{tcolorbox}

% -----------------------------------------------------------------------
%  PIED - PReSENCE DANS LES SUJETS
% -----------------------------------------------------------------------

\begin{tcolorbox}[
  enhanced, breakable,
  colback=colGrayLight,
  colframe=colSeparator,
  left=4pt, right=4pt, top=3pt, bottom=3pt,
  before skip=4pt, after skip=0pt,
  fonttitle=\small\bfseries,
  title={Presence dans les sujets analyses (2013-2026)},
]
\begin{center}
\footnotesize\renewcommand{\arraystretch}{1.15}
\begin{tabular}{lccccccccccccc}
\toprule
\textbf{Annee} & 13 & 14 & 15 & 16 & 17 & 18 & 19 & 20 & 21 & 22 & 23 & 24 & 25 \\
\midrule
E1 (commune)   & - & - & - & - & - & - & - & $\checkmark$ & - & - & - & $\checkmark$ & - \\
E2 discipl. IE & - & - & - & $\checkmark$ & $\checkmark$ & $\checkmark$ & - & $\checkmark$ & $\checkmark$ & - & $\checkmark$ & $\checkmark$ & - \\
\bottomrule
\end{tabular}
\end{center}
\smallskip
\footnotesize\textbf{Sous-themes frequents :}
bilan de puissance $P_a \to P_t \to P_{meca} \to P_u$ (8/13) - vitesse synchrone et glissement (9/13) - couple EM (5/13) - modele Laplace mecanique (2/13) - demarrage etoile-triangle (3/13)

\medskip
\textbf{Ordres de grandeur (reperes rapides).}
\begin{itemize}[leftmargin=1.3em]
  \item Reseau BT FR : tension phase-neutre typique \SI{230}{\volt}, tension phase-phase typique \SI{400}{\volt}, frequence \SI{50}{\hertz}.
  \item Courants d'etude frequents en sujet : de quelques \si{\ampere} a quelques dizaines de \si{\ampere}.
  \item En conversion d'energie, un rendement attendu reste generalement entre 0{,}8 et 0{,}98 selon la technologie.
\end{itemize}

\medskip
\textbf{Synthese express (a refaire sans support).}
\begin{enumerate}[leftmargin=1.4em]
  \item Poser les hypotheses du modele et les conventions de signe.
  \item Ecrire 2 a 4 relations cles avec unites.
  \item Realiser un calcul numerique complet et verifier l'ordre de grandeur.
  \item Justifier le resultat par un argument technique de niveau CAPET.
\end{enumerate}

\end{tcolorbox}

\CAPETENDFILE









